% Canonical derivation source for the exact isotropic ZGV critical ring.
% This file is a LaTeX fragment and is included by the Supplementary Information.

\section{Exact isotropic Rayleigh--Lamb branch and its Morse--Bott ring}
\label{sec:isotropic-rayleigh-lamb-derivation}

We derive the symmetric Rayleigh--Lamb equation from the displacement
potentials, retain the physical traction matrix before any algebraic
division, and then construct the determinant used by the numerical solver.
The plate occupies
\(
  \{(x,y,z):-h\leq z\leq h\}
\)
and is infinite in the in-plane coordinates.  Its density is \(\rho>0\),
and its isotropic stiffness is specified by the Lam\'e moduli \(\lambda\)
and \(\mu\), with
\begin{equation}
  c_L^2=\frac{\lambda+2\mu}{\rho},
  \qquad
  c_T^2=\frac{\mu}{\rho}.
  \label{eq:isotropic-wave-speeds}
\end{equation}
We use the time convention \(\exp[-\mathrm{i}\omega t]\).  Isotropy lets us
rotate any in-plane wave vector onto the \(x\) axis while deriving the scalar
equation; the two-dimensional wave-vector dependence is restored at the end.

\subsection{Helmholtz potentials and symmetric parity}

In the sagittal \((x,z)\) plane, the displacement
\(\mathbf{u}=(u_x,0,u_z)\) obeys Navier's equation
\begin{equation}
  \rho\,\partial_t^2\mathbf{u}
  = (\lambda+\mu)\nabla(\nabla\mathbin{\cdot}\mathbf{u})
    +\mu\nabla^2\mathbf{u}.
  \label{eq:navier-isotropic}
\end{equation}
Introduce a longitudinal potential \(\phi\) and a sagittal shear potential
\(\psi\) by
\begin{equation}
  \mathbf{u}
  =\nabla\phi+\nabla\mathbin{\times}(\psi\,\mathbf{e}_y),
  \qquad
  u_x=\partial_x\phi-\partial_z\psi,
  \qquad
  u_z=\partial_z\phi+\partial_x\psi.
  \label{eq:helmholtz-sagittal}
\end{equation}
The divergence of the curl term is zero and the curl of the gradient term is
zero.  Taking the divergence and the sagittal curl of
Eq.~\eqref{eq:navier-isotropic}, respectively, gives the two scalar wave
equations
\begin{equation}
  \left(\nabla^2-c_L^{-2}\partial_t^2\right)\phi=0,
  \qquad
  \left(\nabla^2-c_T^{-2}\partial_t^2\right)\psi=0.
  \label{eq:potential-wave-equations}
\end{equation}
For a factor \(\exp[\mathrm{i}(kx-\omega t)]\), define the longitudinal and
shear thickness wavenumbers on their principal complex branches by
\begin{equation}
  p^2=\frac{\omega^2}{c_L^2}-k^2,
  \qquad
  s^2=\frac{\omega^2}{c_T^2}-k^2.
  \label{eq:thickness-wavenumbers}
\end{equation}
The symmetric Lamb parity is
\begin{equation}
  u_x(x,-z)=u_x(x,z),
  \qquad
  u_z(x,-z)=-u_z(x,z).
  \label{eq:symmetric-displacement-parity}
\end{equation}
Equation~\eqref{eq:helmholtz-sagittal} therefore requires an even
longitudinal potential and an odd shear potential.  With physical potential
amplitudes \(A\) and \(B\), the complete symmetric separated ansatz is
\begin{equation}
  \phi=A\cos(pz)\,\mathrm{e}^{\mathrm{i}(kx-\omega t)},
  \qquad
  \psi=B\sin(sz)\,\mathrm{e}^{\mathrm{i}(kx-\omega t)}.
  \label{eq:symmetric-potentials}
\end{equation}
Suppressing the common exponential factor from now on, direct substitution
gives
\begin{align}
  u_x(z)
  &=\mathrm{i}kA\cos(pz)-sB\cos(sz),
  \label{eq:symmetric-ux}\\
  u_z(z)
  &=-pA\sin(pz)+\mathrm{i}kB\sin(sz).
  \label{eq:symmetric-uz}
\end{align}
The first expression is even and the second is odd, so the chosen potential
parities reproduce Eq.~\eqref{eq:symmetric-displacement-parity} explicitly.

\subsection{Face tractions and the original determinant}

For an isotropic solid,
\(
  \boldsymbol{\sigma}
  =\lambda(\nabla\mathbin{\cdot}\mathbf{u})\mathbf{I}
   +2\mu\boldsymbol{\varepsilon}
\), where
\(
  \boldsymbol{\varepsilon}
  =(\nabla\mathbf{u}+\nabla\mathbf{u}^{\mathsf T})/2
\).
The derivatives of Eqs.~\eqref{eq:symmetric-ux} and
\eqref{eq:symmetric-uz} give
\begin{align}
  \varepsilon_{xx}
  &=-k^2A\cos(pz)-\mathrm{i}ksB\cos(sz),\\
  \varepsilon_{zz}
  &=-p^2A\cos(pz)+\mathrm{i}ksB\cos(sz),\\
  2\varepsilon_{xz}
  &=-2\mathrm{i}kpA\sin(pz)+(s^2-k^2)B\sin(sz).
  \label{eq:symmetric-strains}
\end{align}
Consequently,
\begin{equation}
  \frac{\sigma_{xz}}{\mu}
  =-2\mathrm{i}kpA\sin(pz)+(s^2-k^2)B\sin(sz).
  \label{eq:shear-traction-through-thickness}
\end{equation}
For the normal stress, the shear-potential terms cancel from the dilatation,
and hence
\begin{align}
  \sigma_{zz}
  &=-\left[\lambda k^2+(\lambda+2\mu)p^2\right]
       A\cos(pz)
    +2\mathrm{i}\mu ksB\cos(sz).                    \label{eq:sigma-zz-first}\\
  \lambda k^2+(\lambda+2\mu)p^2
  &=\rho\omega^2-2\mu k^2
    =\mu\left(\frac{\omega^2}{c_T^2}-2k^2\right)
    =\mu(s^2-k^2),                                    \label{eq:normal-factor-reduction}\\
  \frac{\sigma_{zz}}{\mu}
  &=-(s^2-k^2)A\cos(pz)+2\mathrm{i}ksB\cos(sz).
  \label{eq:normal-traction-through-thickness}
\end{align}
Here Eq.~\eqref{eq:normal-factor-reduction} uses both relations in
Eq.~\eqref{eq:thickness-wavenumbers} and \(\rho c_T^2=\mu\).

Traction freedom means \(\sigma_{xz}=\sigma_{zz}=0\) on each face.  At the
upper face \(z=+h\), Eqs.~\eqref{eq:shear-traction-through-thickness} and
\eqref{eq:normal-traction-through-thickness} become
\begin{equation}
  \underbrace{\begin{pmatrix}
    -2\mathrm{i}kp\sin(ph) & (s^2-k^2)\sin(sh)\\
    -(s^2-k^2)\cos(ph) & 2\mathrm{i}ks\cos(sh)
  \end{pmatrix}}_{\displaystyle T_S(k,\omega)}
  \begin{pmatrix}A\\B\end{pmatrix}=\mathbf{0}.
  \label{eq:traction-plus}
\end{equation}
At the lower face \(z=-h\), oddness of \(\sigma_{xz}\) and evenness of
\(\sigma_{zz}\) give the separate equations
\begin{equation}
  \begin{pmatrix}
    2\mathrm{i}kp\sin(ph) & -(s^2-k^2)\sin(sh)\\
    -(s^2-k^2)\cos(ph) & 2\mathrm{i}ks\cos(sh)
  \end{pmatrix}
  \begin{pmatrix}A\\B\end{pmatrix}=\mathbf{0}.
  \label{eq:traction-minus}
\end{equation}
The first row of Eq.~\eqref{eq:traction-minus} is the negative of the first
row of Eq.~\eqref{eq:traction-plus}, while the second rows coincide.
Therefore either face supplies the same two independent scalar conditions,
and satisfaction on one face enforces satisfaction on the other for a mode
of the stated parity.

Before any amplitude rescaling, the determinant of the physical traction
matrix in Eq.~\eqref{eq:traction-plus} is
\begin{align}
  F_S(k,\omega)
  &=\det T_S                                                        \notag\\
  &=\left[-2\mathrm{i}kp\sin(ph)\right]
    \left[2\mathrm{i}ks\cos(sh)\right]
   -\left[(s^2-k^2)\sin(sh)\right]
    \left[-(s^2-k^2)\cos(ph)\right]                                \notag\\
  &=4k^2ps\cos(sh)\sin(ph)
    +(s^2-k^2)^2\sin(sh)\cos(ph).
  \label{eq:original-traction-determinant}
\end{align}
Thus \(F_S=0\) is the usual symmetric Rayleigh--Lamb equation wherever the
chosen potential basis is nonsingular.

\subsection{Desingularization and the shear cutoff}

At the shear cutoff \(s=0\), the basis function
\(B\sin(sz)\) collapses identically and both entries in the second column of
\(T_S\) vanish.  The resulting factor of \(s\) in
Eq.~\eqref{eq:original-traction-determinant} is therefore a basis
normalization zero, not a physical dispersion condition.  Define the regular
shear amplitude
\begin{equation}
  \widehat B=sB,
  \qquad
  B\sin(sz)=\widehat B\,\frac{\sin(sz)}{s}.
  \label{eq:regular-shear-amplitude}
\end{equation}
For \(s\neq0\), the physical and regular amplitude vectors are related by
\(
  (A,B)^{\mathsf T}=\operatorname{diag}(1,s^{-1})
  (A,\widehat B)^{\mathsf T}
\).  Hence the regular traction matrix is
\begin{align}
  \widetilde T_S
  &=T_S\operatorname{diag}(1,s^{-1}) \notag\\
  &=\begin{pmatrix}
    -2\mathrm{i}kp\sin(ph)
      &(s^2-k^2)\dfrac{\sin(sh)}{s}\\[6pt]
    -(s^2-k^2)\cos(ph)
      &2\mathrm{i}k\cos(sh)
  \end{pmatrix}.
  \label{eq:regular-traction-matrix}
\end{align}
Taking its determinant entry by entry gives
\begin{align}
  D_S(k,\omega)
  :=\det\widetilde T_S
  &=(s^2-k^2)^2\frac{\sin(sh)}{s}\cos(ph)
    +4k^2p\cos(sh)\sin(ph).
  \label{eq:desingularized-determinant}
\end{align}
For \(s\neq0\), Eqs.~\eqref{eq:original-traction-determinant} and
\eqref{eq:desingularized-determinant} satisfy exactly
\begin{equation}
  F_S=sD_S.
  \label{eq:original-regular-determinant-relation}
\end{equation}
They therefore have identical roots away from the shear cutoff.  The regular
matrix also defines the correct analytic continuation through that cutoff.

Indeed, the even entire function multiplying the second potential has the
Taylor series
\begin{equation}
  \frac{\sin(sh)}{s}
  =h-\frac{s^2h^3}{3!}+\frac{s^4h^5}{5!}
    -\frac{s^6h^7}{7!}+O(s^8),
  \qquad s\longrightarrow0.
  \label{eq:sine-ratio-series}
\end{equation}
Together with \(\cos(sh)=1-s^2h^2/2+O(s^4)\), this gives the finite matrix and
determinant limits
\begin{align}
  \widetilde T_S\big|_{s=0}
  &=\begin{pmatrix}
    -2\mathrm{i}kp\sin(ph)&-k^2h\\
    k^2\cos(ph)&2\mathrm{i}k
  \end{pmatrix},                                                    \label{eq:regular-matrix-s-cutoff}\\
  D_S\big|_{s=0}
  &=k^4h\cos(ph)+4k^2p\sin(ph).                                    \label{eq:regular-determinant-s-cutoff}
\end{align}
At this cutoff \(\omega=c_Tk\) and
\(
  p^2=(c_T^2/c_L^2-1)k^2
\); Eq.~\eqref{eq:regular-determinant-s-cutoff} is not identically zero.
It retains a zero only when the two regular face-traction equations are
linearly dependent, which is the physical condition.

The other square root causes no loss of analyticity: \(\cos(ph)\) and
\(p\sin(ph)\) are entire functions of \(p^2\), while
\(\cos(sh)\) and \(\sin(sh)/s\) are entire functions of \(s^2\).
Equation~\eqref{eq:thickness-wavenumbers} therefore supplies analytic
continuation into evanescent regions.  For real \(k\) and \(\omega\), the
continued determinant remains real even when \(p\) or \(s\) is purely
imaginary.

\subsection{Implicit dispersion derivatives}

Let \((k_0,\omega_0)\) be a root of \(D_S\) with
\begin{equation}
  D_{S,\omega}(k_0,\omega_0)\neq0.
  \label{eq:nonzero-frequency-derivative}
\end{equation}
The implicit-function theorem then supplies a unique smooth local branch
\(\omega=\omega(k)\) satisfying
\begin{equation}
  D_S(k,\omega(k))=0.
  \label{eq:implicit-dispersion-branch}
\end{equation}
In the following equations, every partial derivative of \(D_S\) is evaluated
at \((k,\omega(k))\).  One total derivative of
Eq.~\eqref{eq:implicit-dispersion-branch} gives
\begin{align}
  0
  &=\frac{\mathrm{d}}{\mathrm{d}k}D_S(k,\omega(k))
    =D_{S,k}+D_{S,\omega}\omega_k,                                  \label{eq:implicit-first-chain}\\
  \omega_k
  &=-\frac{D_{S,k}}{D_{S,\omega}}.                                  \label{eq:implicit-first-derivative}
\end{align}
Differentiating Eq.~\eqref{eq:implicit-first-chain} once more, including the
dependence of each partial derivative on both arguments, yields
\begin{align}
  0
  &=\frac{\mathrm{d}}{\mathrm{d}k}
    \left(D_{S,k}+D_{S,\omega}\omega_k\right)                         \notag\\
  &=\left(D_{S,kk}+D_{S,k\omega}\omega_k\right)
   +\left(D_{S,\omega k}+D_{S,\omega\omega}\omega_k\right)\omega_k
   +D_{S,\omega}\omega_{kk}                                         \notag\\
  &=D_{S,kk}+2D_{S,k\omega}\omega_k
    +D_{S,\omega\omega}\omega_k^2+D_{S,\omega}\omega_{kk},
  \label{eq:implicit-second-chain}
\end{align}
where equality of mixed derivatives was used in the last line.  Solving this
equation gives the general second derivative
\begin{equation}
  \omega_{kk}
  =-\frac{D_{S,kk}+2D_{S,k\omega}\omega_k
          +D_{S,\omega\omega}\omega_k^2}
         {D_{S,\omega}}.
  \label{eq:implicit-second-derivative}
\end{equation}
A zero-group-velocity point satisfies \(\omega_k(k_0)=0\).  Under
Eq.~\eqref{eq:nonzero-frequency-derivative}, this is equivalent to
\(D_{S,k}(k_0,\omega_0)=0\).  Substitution into
Eq.~\eqref{eq:implicit-second-derivative} leaves the radial curvature
\begin{equation}
  a:=\omega_{kk}(k_0)
  =-\left.\frac{D_{S,kk}}{D_{S,\omega}}
    \right|_{(k_0,\omega_0)}.
  \label{eq:zgv-curvature}
\end{equation}
Thus, with \(q=k-k_0\), Taylor's theorem gives
\begin{equation}
  \omega(k)=\omega_0+\frac{a}{2}q^2+O(q^3).
  \label{eq:isotropic-radial-normal-form}
\end{equation}

For the selected first symmetric branch (the \(S_1\)-ZGV branch within the
configured search window) and the nondimensional reference values
\(
  h=\rho=\mu=1
\) and \(\lambda=2\), so that \(c_T=1\) and \(c_L=2\), the independent
arbitrary-precision refinement of \(D_S=D_{S,k}=0\), seeded only by the
float-valued branch search, is
\begin{equation}
  k_0=0.8042173193715181,
  \qquad
  \omega_0=2.8517587749600901,
  \qquad
  a=1.1968627250739301>0.
  \label{eq:reference-isotropic-zgv-values}
\end{equation}
The executable check associated with this derivation differentiates
Eq.~\eqref{eq:desingularized-determinant} at arbitrary precision, verifies
\(D_S=D_{S,k}=0\), verifies \(D_{S,\omega}\neq0\), and compares all three
values in Eq.~\eqref{eq:reference-isotropic-zgv-values} with both the Task-4
solver and the generated reference artifact.  This is a separate
high-precision refinement of the same exact scalar determinant; the
independent thickness-resolved spectral cross-check is supplied by the
numerical validation rather than claimed here.

\subsection{Cartesian Hessian and the Morse--Bott ring}

Restore the two-dimensional wave vector
\(
  \mathbf{k}=(k_x,k_y)
\), and write
\begin{equation}
  r=|\mathbf{k}|=(k_x^2+k_y^2)^{1/2},
  \qquad
  \mathbf{e}_r=\frac{\mathbf{k}}{r},
  \qquad
  \mathbf{e}_\theta=(-e_{r,y},e_{r,x}).
  \label{eq:polar-wavevector-basis}
\end{equation}
For \(r>0\), isotropy makes the selected branch radial:
\(
  \omega^{(0)}(\mathbf{k})=f(r)
\), where \(f(k)=\omega(k)\) in the one-dimensional derivation above.  Since
\begin{equation}
  \partial_{k_i}r=\frac{k_i}{r},
  \qquad
  \partial_{k_j}\left(\frac{k_i}{r}\right)
  =\frac{\delta_{ij}}{r}-\frac{k_i k_j}{r^3},
  \label{eq:radial-coordinate-derivatives}
\end{equation}
the Cartesian gradient and Hessian follow component by component:
\begin{align}
  \partial_{k_i}\omega^{(0)}
  &=f'(r)\frac{k_i}{r},                                               \label{eq:radial-cartesian-gradient}\\
  \partial_{k_j}\partial_{k_i}\omega^{(0)}
  &=f''(r)\frac{k_i k_j}{r^2}
    +f'(r)\left(\frac{\delta_{ij}}{r}
                 -\frac{k_i k_j}{r^3}\right).                        \label{eq:radial-cartesian-hessian-components}
\end{align}
In dyadic form this is
\begin{equation}
  \nabla_{\mathbf{k}}^2\omega^{(0)}
  =f''(r)\,\mathbf{e}_r\mathbf{e}_r^{\mathsf T}
   +\frac{f'(r)}{r}
    \left(\mathbf{I}-\mathbf{e}_r\mathbf{e}_r^{\mathsf T}\right)
  =f''(r)\,\mathbf{e}_r\mathbf{e}_r^{\mathsf T}
   +\frac{f'(r)}{r}\,\mathbf{e}_\theta\mathbf{e}_\theta^{\mathsf T}.
  \label{eq:radial-cartesian-hessian}
\end{equation}
Consequently the orthogonal polar basis diagonalizes the full Cartesian
Hessian:
\begin{equation}
  \begin{pmatrix}\mathbf{e}_r&\mathbf{e}_\theta\end{pmatrix}^{\mathsf T}
  \nabla_{\mathbf{k}}^2\omega^{(0)}
  \begin{pmatrix}\mathbf{e}_r&\mathbf{e}_\theta\end{pmatrix}
  =\begin{pmatrix}
    f''(r)&0\\[2pt]
    0&f'(r)/r
  \end{pmatrix}.
  \label{eq:radial-hessian-diagonalization}
\end{equation}
At \(r=k_0\), the ZGV conditions \(f'(k_0)=0\) and \(f''(k_0)=a\) reduce
Eq.~\eqref{eq:radial-hessian-diagonalization} to
\begin{equation}
  \nabla_{\mathbf{k}}^2\omega^{(0)}\big|_{r=k_0}
  =a\,\mathbf{e}_r\mathbf{e}_r^{\mathsf T},
  \qquad
  \operatorname{spec}
  \left(\nabla_{\mathbf{k}}^2\omega^{(0)}\big|_{r=k_0}\right)
  =\{a,0\}.
  \label{eq:zgv-ring-hessian}
\end{equation}

We now state precisely when the rotationally generated ZGV circle is a
Morse--Bott critical manifold.  The hypotheses separate the local scalar
geometry from the spectral condition that makes a single elastic branch
well defined.

\begin{theorem}[Isotropic ZGV Morse--Bott ring]
\label{thm:morse-bott-ring}
Let a lossless homogeneous isotropic plate have a symmetric guided-wave
eigenbranch in an open annulus about \(r=k_0>0\).  Assume:
\begin{enumerate}
  \item the eigenvalue is simple throughout that annulus and is separated
        from every other Lamb and shear-horizontal eigenvalue by a uniform
        positive mass-normalized eigengap;
  \item the branch is \(C^3\) in \(\mathbf{k}\) and rotational invariance
        gives \(\omega^{(0)}(\mathbf{k})=f(|\mathbf{k}|)\);
  \item the regular determinant satisfies
        \(D_S(k_0,\omega_0)=D_{S,k}(k_0,\omega_0)=0\) and
        \(D_{S,\omega}(k_0,\omega_0)\neq0\); and
  \item the radial curvature
        \(a=-D_{S,kk}/D_{S,\omega}\) at \((k_0,\omega_0)\) is nonzero.
\end{enumerate}
Then, in a sufficiently small annular neighbourhood, the critical set of
\(\omega^{(0)}\) is exactly
\begin{equation}
  \Gamma_0=\{\mathbf{k}\in\mathbb{R}^2:|\mathbf{k}|=k_0\}.
  \label{eq:isotropic-critical-ring}
\end{equation}
Moreover, at every \(\mathbf{k}\in\Gamma_0\),
\begin{equation}
  \ker\nabla_{\mathbf{k}}^2\omega^{(0)}
  =\operatorname{span}\{\mathbf{e}_\theta\}
  =T_{\mathbf{k}}\Gamma_0,
  \qquad
  \left.\nabla_{\mathbf{k}}^2\omega^{(0)}\right|_{N_{\mathbf{k}}\Gamma_0}
  =[a].
  \label{eq:morse-bott-kernel-and-normal-form}
\end{equation}
Thus \(\Gamma_0\) is a one-dimensional nondegenerate Morse--Bott critical
manifold.  If \(a>0\), it is a radial minimum manifold; if \(a<0\), it is a
radial maximum manifold.
\end{theorem}

\paragraph{Proof.}
The nonzero value of \(D_{S,\omega}\) gives a unique local scalar branch by
the implicit-function theorem, and the simple uniformly gapped elastic
eigenvalue identifies that scalar root with one smooth physical mode.  From
Eq.~\eqref{eq:implicit-first-derivative},
\(f'(k_0)=0\), while Eq.~\eqref{eq:zgv-curvature} gives
\(f''(k_0)=a\neq0\).  Continuity of \(f''\) permits a sufficiently small
\(\delta>0\) such that \(f''(r)\) has the sign of \(a\) whenever
\(|r-k_0|<\delta\).  Equivalently, the fundamental theorem of calculus gives
\begin{equation}
  f'(r)=f'(k_0)+\int_{k_0}^{r}f''(\xi)\,\mathrm{d}\xi
       =(r-k_0)\int_0^1 f''\bigl(k_0+\tau(r-k_0)\bigr)\,\mathrm{d}\tau.
  \label{eq:local-critical-radius-factorization}
\end{equation}
The integral factor is nonzero in that interval, so \(f'(r)=0\) there if and
only if \(r=k_0\).  Equation~\eqref{eq:radial-cartesian-gradient} then proves
that the critical set in the annulus is exactly \(\Gamma_0\).

For each point on the ring,
Eq.~\eqref{eq:zgv-ring-hessian} maps \(\mathbf{e}_r\) to
\(a\mathbf{e}_r\) and maps \(\mathbf{e}_\theta\) to zero.  Since \(a\neq0\),
its kernel is precisely \(\operatorname{span}\{\mathbf{e}_\theta\}\), which
is the tangent line to the circle, while the restriction to the radial
normal line is the nonsingular one-by-one form \([a]\).  These are exactly
the critical-set and normal-nondegeneracy conditions in the definition of a
Morse--Bott function.  The sign of \(a\) supplies the final minimum/maximum
classification. \(\square\)

For the reference branch in Eq.~\eqref{eq:reference-isotropic-zgv-values},
\(a>0\), so the normal Morse index is zero.  The tangential zero eigenvalue is
not an accidental higher-order radial degeneracy: it is enforced by the
continuous rotational symmetry.  The cases \(a=0\),
\(D_{S,\omega}=0\), or a closing modal eigengap violate distinct hypotheses
of Theorem~\ref{thm:morse-bott-ring} and require a higher-order or
matrix-valued local reduction.
